% Chapter VIII: Geometrically Contradictory Images
% Source pages: 168–189

\chapter{Geometrically Contradictory Images}

\markboth{Chapter VIII. Geometrically Contradictory Images}{Chapter VIII}

\srcpage{168}

In Chapter~III and the corresponding Appendices it was shown that the representation of
nearby regions of space without any ``errors'' is in general impossible.  However, the
conclusions reached there do not exhaust the question of the fundamental possibility of
an ideal mapping of space onto a plane.  Let us examine this question again from
completely different standpoints, making no distinction between near and far space.

In the depiction of perceptual space on a plane, the problem of the ambiguity of the
representation arises: if each point of the pictorial space corresponds to one point in
the picture plane, the converse does not hold---one point of the picture may correspond
to an infinite set of points of the pictorial space.\footnote{Below, wherever the
  ambiguity of the mapping is discussed, it is always understood that to one point of
  the picture there corresponds an infinite set of points of the pictorial space.}

To explain the essence of this ambiguity, let us return to Ill.~14.  The point~$a$,
defined by the upper right corner of the image of cube~$A$, may correspond to different
points of the space that the picture maps.  As is clear from diagrams~$B$ and~$C$, it
may, for example, belong either to the front or to the rear face of a cube situated in
three-dimensional space.  Consequently, the position of point~$a$ in the picture does
not yet determine the position of the corresponding point in space (though the converse
may hold: the position of a point in space may uniquely determine its position in the
picture plane; this depends on the method of establishing the correspondence, and in
particular ordinary linear perspective has this property).  Perhaps the most natural
explanation of the arising ambiguity is the following: all points of three-dimensional
space lying on one line of sight fall on one point of the retina and are in this sense
indistinguishable.  They may obscure one another (if opaque material bodies are located
there), but they cannot fall on different points of the retina.  Consequently, each
point of the retina (and hence each corresponding point of the picture) maps an infinite
set of points lying on the corresponding line of sight.

The ambiguous correspondence of points of the picture plane to points of pictorial space,
described here, prompts artists to employ special measures for resolving this ambiguity.
In Fig.~14, diagrams~$B$ and~$C$ already provide the required unambiguity in the
depiction of a cube---this is achieved by using the principal depth cue, occlusion (the
faces of the cube closer to the viewer occlude the more distant ones).

\srcpage{169}

There are known cases in which artists deliberately violated the ``rules'' for resolving
the ambiguity of an image, thereby creating ``impossible pictures.''  Usually this had
the character of a jest; in particular Hogarth left us drawings of this
kind.\footnote{See, for example: R.~Gregory.  \emph{Razumny glaz} [The Intelligent Eye].
  Moscow, 1972, p.~58.}  In recent years the well-known Dutch graphic artist Escher
(sometimes written Esher) devoted much attention to creating ``strange'' drawings, in
which he combined geometry with drawing with great inventiveness, treating the resulting
strangeness with complete seriousness.\footnote{``The Graphic Work of M.~C.~Escher.''
  Hawthorn Books, Inc., publishers.  New York.  Two of his engravings are reproduced in
  R.~Gregory's book.}  One of his engravings appears to depict something like a
two-storey gazebo with people inside it dressed in antique European costumes (Ill.~75).
The engraving has certain artistic merits and creates a feeling of somewhat melancholy
pensiveness.

More attentive examination of the engraving may throw into confusion a person who loves
clarity and unambiguity in everything: the nobleman and the lady looking into the distance
and standing on different floors between the same columns are not looking in the same
direction but in mutually perpendicular directions; the people climbing a portable ladder
to the upper floor have placed it in a way in which it cannot stand (it stands inside
the gazebo but leans against it from outside); the bases and capitals of a number of
columns are not on the same vertical; and so on.  If one tries to understand the nature
of the errors in the drawing, it is easy to see in each case that Escher deliberately
violates obvious rules of drawing: since one column can occlude another situated in the
depth of the pictorial space, it suffices to draw the capital of one and the base of
another as belonging to the same structure, and an ``impossible'' column comes into being.
This simple device can be carried out quite inconspicuously for the viewer, since by
reason of the ambiguity of the mapping of space onto the picture discussed above, both
``correct'' columns would occupy one and the same position in the picture plane.

The example given here allows one to illuminate one feature of the psychology of
perception.  Since the given retinal image (which in the end reduces to a set of coloured
patches, lines, and so on) may, by reason of the ambiguity under discussion, correspond
to an infinite set of different spatial formations, the human perceptual system selects
from this infinite set of interpretations of the retinal image the single most plausible
one, corresponding to life experience.  This selection occurs in the great majority of
cases infallibly and, most importantly, subconsciously.  Helmholtz called this
``unconscious inference.''  It is important to note here that only a small fraction of
the information that can be extracted from the retinal image reaches consciousness, but
it is the vitally important part.  In the pro-

\srcpage{170}

\figblock{75}{M.~C.~Escher.  ``Belvedere.''  Lithograph.}

cess of perception there occurs a ``filtering'' of everything that is secondary in the
given situation.  For example, when recognising a friend, no one does this through
conscious analysis of the coloured patches and lines of the retinal image, but ``simply''
recognises---the process of recognition occurs at some lower level than conscious
analysis.  In the recognition process important roles are played by the features of the
face, the hairstyle, hair colour, and so on, but no one notices the position of each
individual hair in the coiffure, even though the retinal image contains this information.

In the process of recognition, from the whole mass of information contained in the retinal
image, exactly as much is extracted as is needed for recognition, and all the rest is
``filtered.''  Human perception, as it were sorting through the possible variants of
interpretation of what is seen, ceases to extract additional information from the
observed collection of coloured patches and lines as soon as the interpretation of the
meaning of these patches and lines becomes unambiguous, so as not to ``clog'' consciousness
with superfluous details.

The recognition process characterised here in the most general terms is very brief and
very economical, which is biologically important.  In life, a person deals with moving
objects, and therefore only a fraction of a second is allotted for the assessment of the
situation---for what was called ``recognition'' above---in order that during the following
fraction of a second the two successive pictures of the surrounding world may be
subconsciously compared, and moving objects distinguished from stationary ones by such
comparison, and so on.  Therefore the ``filtering'' of the superfluous contained in the
retinal image is a strict biological necessity that took shape in the process of evolution.

All that has been said above applies to pictures as well.  When looking at a painted work,
a person extracts from the mass of information contained in the picture exactly as much
as is needed for his perception to give an unambiguous answer to the question ``what is
this?''---though here this is not a biological necessity but a consequence of perceptual
laws already formed for other purposes.  After ``recognition'' the corresponding emotions
may arise, and so on.  To compel oneself to extract additional information from the
picture requires a volitional effort, which we often call ``looking more carefully''; in
this process a more conscious analysis of the details of the picture begins.  The next,
still higher level on this path is the logical and aesthetic analysis of the compatibility
of the picture's details, in particular the internal self-consistency of what is depicted.

In the majority of cases the ordinary viewer confines himself to the first of the levels
of perception named here, since, as has been noted, ``recognition'' is sufficient for a
person to move from the realm of ``unconscious inferences'' to the realm of emotional
perception.  He may admire the picture, emotions may be aroused in him that the artist
wished to arouse in creating his work, and so on.  How serious a barrier our psyche erects
against the penetration into consciousness of ``superfluous'' information is shown by the
type of drawing---familiar to all---created for testing powers of observation.  In
magazines one often encounters a drawing in which the reader is invited to find, say,
15 errors.  Although all 15 errors are immediately ``imprinted'' on the retina, the
reader often has to spend much time and effort to find them.

\srcpage{171}

The reason is that these errors are distributed in the drawing so as not to impede the
first, principal level of perception, the process of ``recognition,'' since once an
unambiguous interpretation of what is drawn has been achieved, the receipt of additional
information is ``cut off,'' and this additional information has to be obtained by applying
a volitional effort.

Even such gross errors in a drawing as those characteristic of the Escher engraving
cited are not immediately conspicuous.  The majority of ordinary viewers to whom this
engraving was shown responded to it emotionally (``I like it,'' ``a beautiful engraving,''
``nothing special,'' and so on) but did not point out the gross errors made by the artist.

The considerations presented above, and in particular Escher's engraving, are capable of
giving rise to reflections about the extent to which ``correct'' resolution of ambiguity
in a picture is generally necessary, and whether the artist is not entitled sometimes to
make ``errors'' deliberately.  Would not ``errors'' in some cases increase the
expressiveness of the picture as a whole?  For after all the abundance of obvious errors
in the engraving does not spoil the general impression it creates, provided one does not
attach too much importance to the photographic accuracy of what is depicted.

One of the differences between medieval painting and the art of the Renaissance is the
following.  If the new art is characterised by the desire to resolve all questions
connected with the unavoidable ambiguity of the space mapped in the picture in such a
way that it contains no contradictions even under scrutiny and analysis, medieval and
more ancient art does not attach decisive importance to this---and moreover: by
deliberately violating these ``photographic'' rules in certain cases, it increases the
overall expressiveness of the picture.  This is yet another manifestation of the
artist's freedom, described above in connection with perspective effects and the use of
drawing methods, and which, as a rule, is conditioned by artistic tasks, even at the
expense of accurate rendering of the visible picture of the external world.

\srcpage{172}

\figblock{76}{Panel of the icon ``Life-giving Spring of the Mother of God,'' 1675.
  Museum of Architecture and Applied Arts (MiAR).}

The author of the icon ``Life-giving Spring of the Mother of God'' (Ill.~76) employs a
device close to the one shown in Escher's engraving: in the panels narrating the miracles
at this spring, he shifts the bases of the corner columns to the foreground and thereby
achieves two artistic effects: these columns separate the panels from each other and,
moreover, as it were enclose the image of the spring within the interior space of the
building.  The geometric contradiction here is not accidental but introduced deliberately.

By geometric contradiction in an image, one by no means understands geometric
``impossibility'' of what is depicted.  Indeed, columns with bases in the foreground and
capitals in the background are geometrically possible---they would be columns inclined
at a large angle to the vertical.  However, such columns do not exist, and therefore a
person looking at the picture will find such a column ``strange,'' yet will nonetheless
feel it to be vertical and know that it is vertical.  By geometric contradiction is
understood the contradiction between the geometric image that arises in the viewer's
perception of the picture and the formal geometric properties that the representation
actually possesses.  Below, the concept of geometric contradiction is to be understood
only in this sense.

The icon-painter used the shift of a column base ``forward'' to obtain the required
artistic effect.  In other cases a shift ``backward'' may be needed.  In one of the
Byzantine eleventh-century miniatures there is a depiction of an elder seated under a
baldachin (Ill.~77).  The lower ends of the posts supporting the baldachin are shifted
backward, so that the front posts do not prevent the gesticulating elder from being shown
in a proper manner.  Thus here too the violation of geometric self-consistency enhances
the artistic effect.

\srcpage{173}

\figblock{77}{Detail of a miniature from the \emph{Ladder} of John Climacus
  (traced drawing).}

Sometimes (though rarely) columns are shifted so as to strengthen the ``effect of
occlusion.''  For example, in the panel of the icon ``Paraskeva Friday'' (Ill.~78) both
the front and rear columns supporting the ceiling have their bases in the foreground, as
a result of which Paraskeva is separated from the viewer by the columns, one of which
even partially occludes the saint.  This unusual manner of depiction can in the given
case be readily explained, since the artist needed in the said panel to show the prison
in which Paraskeva was confined.  In both examples the violation of geometric
self-consistency enhances the content-artistic effect.  These two examples show that the
medieval artist, when introducing geometrically contradictory elements into his image,
understood perfectly well why it was needed in each particular case.

As already mentioned in Chapter~IV, the Old Russian artist, compelled by the objective
properties of perceptual perspective to partially distort the visible picture, extended
this ``right'' of his and began to alter the form of objects in the picture (and hence
in the imagined real space) when this seemed to him artistically justified.  The same is
shown by the techniques of Old Russian art examined in the present chapter.  If violating
the geometric self-consistency of an image while preserving the position of the vertical
axes of the columns in the picture could produce a positive artistic effect, the artist
was not afraid to shift the axes themselves in the horizontal direction in pursuit of his
artistic aim.  With such shifts in the picture plane, the resulting effect was no longer
directly connected with the ambiguity of the inverse mapping of the depicted onto external
space, but effectively amounted to a distortion of the external forms of objects in the
imagined real space.  Violation of the geometric self-consistency of the external space
reconstructed from the picture in this case effectively reduced to the appearance of
``impossible'' constructions, each individual element of which is reasonable but which
as a whole are structurally meaningless.

\srcpage{174}

\figblock{78}{Panel of the icon ``Paraskeva Friday,'' sixteenth century.
  Museum of Architecture and Applied Arts (MiAR).}

To explain this, let us turn to the architectural background of the icon
``Presentation in the Temple,'' shown in Ill.~28.  The four-column construction shown
in the upper left contains the type of distortions discussed above, while the analogous
construction in the upper right of the icon is of a completely different character.
The arrangement of columns there is structurally meaningless and could only have arisen
through horizontal shifts of the columns in the picture.  What guided the artist when
introducing these distortions?  In the left image he was constrained by the narrow strip
at the top of the architectural structure, in which the bases of the columns had to be
placed, while in the right image it seemed aesthetically justified to space the columns
at approximately equal intervals from one another and to limit their number to four.

\srcpage{175}

For the question examined in the present chapter, what matters is something else: why
such distortions are permissible at all.  The answer has already been given above: since
the inverse mapping of what is depicted onto real space is ambiguous, a person in the
process of visual perception subconsciously first sorts through possible perceptual
hypotheses about the meaning of what is depicted, using only the minimum information
sufficient for an unambiguous resolution of the cardinal question ``what is this?''
Therefore the combination of ``roof plus four columns'' immediately makes the meaning of
what is depicted clear, and the viewer subconsciously loses interest in the recognised
details of the painted work, the more so since these details belong to the periphery of
the image and in essence should not attract his particular attention, which should be
concentrated on the lower part of the picture.

Analogous considerations led to the artist not depicting the fourth column supporting
the roof above the head of Mary in the second plane of the picture.  He probably found
it aesthetically wrong if a ``palisade'' of columns were to arise visually between Mary
and the angel appearing to her; he understood that the clear configuration of the roof
and the three clearly visible columns would give the viewer the sense that ``somewhere
behind, there is a fourth column too,'' i.e.\ its absence would not impede ``recognition.''

Thus, in the cases examined there are present (1)~the type of shifts already discussed
above, along the line of sight without change in the position of the vertical axis of the
transformed column (and let us include in this type the ``disappearance'' of a column),
and (2)~explicit horizontal shifts of columns.

The appearance of ``impossible'' constructions in icon painting as a result of horizontal
shifts of columns is not infrequent.  For example, in compositions on the theme of the
``Eucharist'' (Communion of the Apostles) in the great majority of cases the posts of
the ciborium are disposed (including by horizontal shifts) so as not to ``interfere
visually'' with Christ giving communion to the apostles, even though in real space the
posts in that case would be unable to support the canopy exactly over the altar.  The
panel of the Royal Doors ``Eucharist'' (end of the fifteenth century) (Ill.~79) is a
typical example of shifts both along the line of sight and in lateral directions.  Here
three columns supporting the ciborium are arranged in a completely natural way, but the
hands of Christ standing behind the altar are shown closer than the front columns, i.e.,
in front of the altar.  This is geometrically impossible---Christ is not shown bending
forward---but by this geometrically contradictory method it proved possible to show
Christ's figure without intersection by the images of the columns.  Here it is not the
columns that have been shifted, but the image of Christ's hands along the lines of sight.
The left rear column is shifted to the right relative to the corresponding corner of the
altar and its image is truncated.  This is entirely justified compositionally, since it
frees the space before Christ from distracting elements of secondary constructions.

\srcpage{176}

One should not think that the problems connected with geometric ambiguity in the mapping
of space onto a plane are limited to the task of depicting columns.  Thus in one of the
compositions of the ``Baptism'' (Ill.~35) the mutual arrangement of the figures of Christ
and John is geometrically contradictory.  In order to depict Christ standing at full
height and John bent over, as tradition required, the artist placed the latter far
``behind'' Christ (or is he to be considered ``hanging'' in the air?).  This composition
developed from another, earlier one in which Christ was depicted waist-deep in the water
and the bent posture of John was justified.  Probably later the depiction of Christ's
full figure became ideologically and artistically more expedient, and the composition
shown in Ill.~35 relies on the conventions already established that are characteristic
of geometrically contradictory images.  The situation is essentially analogous in the
depiction of the Evangelist Mark (Ill.~55): the table at which he is writing is shifted
along the line of sight and is placed not in front of the Evangelist but to his left.
Yet this composition made it possible to show the full figure of Mark with throne and
footrest, not in the least occluded by the table.

On the Pskov icon of the fourteenth century ``Deesis with Saints'' (Ill.~80), the Mother
of God, the Forerunner, Barbara, and Paraskeva are by the meaning of the scene before
Christ, yet in the image they are ``shifted'' so that Christ does not end up in the depth
of the depicted space.\footnote{Christ's primacy is also emphasised by other means: his
  figure is hierarchically enlarged, and those standing left and right are shown ``one
  above the other.''  This allows the area of the icon allotted to the depiction of Christ
  to be expanded and compositionally justifies the exaggerated height of his figure.}
To the same end, the front legs of Christ's throne almost touch the lower boundary of
the image, and the footrest is ``shifted'' under the throne, which is geometrically
absurd but also allows the figure of Christ to be brought closer to the viewer.  The
contours of the throne occlude the figures of the Mother of God and the Forerunner, which
is impossible given the shown position of their feet.  All this geometric contradiction
has one aim: to bring the image of Christ as close as possible to those contemplating
the icon by every available means, to emphasise his supremacy.

\srcpage{177}

\figblock{79}{``Eucharist.''  Panel of the Royal Doors, late fifteenth century.
  State Russian Museum (GRM).}

\figblock{80}{``Deesis'' with Saints.  Second half of the fourteenth century.
  Novgorod State Museum (NGM).}

To the type of geometrically contradictory images under consideration one must also
include those in which, in depicting for example human crowds or large military masses,
the number of heads and the number of legs shown do not agree.  Thus in the Pskov
sixteenth-century icon ``Nativity,''\footnote{A.~Ovchinnikov, N.~Kishilov.
  \emph{Zhivopis' drevnego Pskova} [Painting of Ancient Pskov].  Moscow, 1976, ill.~48.}
the three Magi hastening to the cradle of Christ are galloping not on three horses but
on something ``incomprehensible'' having one head, two tails, four legs.  Yet we
subconsciously picture a group of three horses in a complex mutual occlusion (though, as
analysis shows, in a geometrically absurd occlusion).

Finally, entirely unexpected types of geometrically contradictory images are possible.
Thus in the Pskov fourteenth-century icon ``Miracle of George and the Dragon,'' the gait
of the horse is emphasised by the depiction of an extra ``knee'' at each of the horse's
front legs.\footnote{See: S.~Yamshchikov.  \emph{Drevnerusskaya zhivopis'.  Novye
  otkrytiya} [Old Russian Painting.  New Discoveries].  Leningrad, 1969, ill.~7.}

The use of geometrically contradictory images is often explained by the semantic, sign
character of medieval painting.  Indeed, if the human visual perception system requires
the ``disconnection'' of superfluous information, and this allowed the icon-painter
(Ill.~28) to depict the upper parts of architectural structures according to the scheme
``roof plus four columns'' without a geometrically correct arrangement of the structural
elements, and to show only three of the four columns supporting the roof above Mary's
head, then one might suppose that the icon presents not true depictions but only their
symbols, signs, simplified relative to their ``correct'' depictions.

\srcpage{178}

One should, however, guard against over-enthusiasm for such interpretations of the
geometric features of Old Russian painting.  If one examines the majority of distortions
similar to those described above, a not unimportant circumstance strikes the eye: almost
all these distortions are concentrated on those secondary details of the image that may
be considered to belong to those ``filtered'' by the visual perception system, i.e.,
inconspicuous or barely noticeable to the viewer.  At the same time, a sign is not at
all obliged to have this property---rather, it should catch the eye.

What has been said does not at all mean that Old Russian painting lacked the sign
character.  It suffices to look at the four-column structure in which Paraskeva is
enclosed (Ill.~78): clearly it cannot serve as a prison---it is simply a sign for a
building.  Indeed, the construction shown belongs to one of the few types of depictions
of buildings ``in general'' used in Old Russian art; its meaning becomes clear only within
the system of the picture as a whole.

The building sign taken as an example does not have the property of ``inconspicuousness''---
it is a conspicuous, conventional, and therefore universally intelligible depiction of any
building (apart from obvious exceptions, e.g.\ a church).  Such signs are highly
constant---one might say standardised---and moreover do not necessarily contain
geometrically contradictory elements.  In contrast to them, the examples of geometrically
contradictory images examined above are interesting in that, like the ``errors'' of
Escher's engraving, their ``incorrectnesses'' also do not catch the eye, are barely
noticeable, and have no sign character.  The ciborium on the Royal Doors in the
``Eucharist'' (Ill.~79) with one ``truncated'' column would not change its meaning if
all four columns were shown complete.  It therefore seems reasonable not to attribute
sign character to images containing geometrically contradictory elements on this ground
alone.  One should understand that the sign criterion belongs to the semantic
classification, while the contradiction discussed in the present chapter belongs to the
classification of the geometric properties of the image.

\srcpage{179}

Returning to the question of the motivating factors that attracted artists' attention to
geometrically contradictory images, let us recall that the examples cited above showed
how abandoning formal adherence to the principles of geometrically self-consistent and
strictly documentary depiction can prove useful and strengthen the artistic impact on
the viewer.  A more complete impression of this effect can be obtained from a comparison
of pictures similar in subject-matter executed in different manners: in one case with
strictly self-consistent resolution of geometric ambiguity, in the other with deliberate
violation of this geometric self-consistency.

Let us compare the two types of depictions of ancient Egyptian married couples, already
discussed in Chapter~I, one of which is characteristic of the Old Kingdom and the other
of the New Kingdom (Ill.~8).  As already noted in Chapter~I, in the Old Kingdom the
dominant role of the husband is expressed not only by his being depicted in front of his
wife, but also by his occluding her.  As the analysis presented there shows, the wife
thereby adopts an impossible pose; consequently in the New Kingdom a transition was made
to a geometrically correct depiction of mutual occlusions.  This change in types of
depictions of married couples is usually interpreted positively, as a victory of realism.

\srcpage{180}

A contrary point of view is, however, also legitimate.  Comparing both compositions of
married couples, one easily notes that the formal geometric absurdity of the more ancient
one does not immediately catch the eye.  To understand it, attentive and thoughtful
analysis of the peculiarities of the depiction is necessary; the ordinary viewer simply
does not notice the unreality of the poses shown in the more ancient image.  But the type
of mutual occlusions is visible immediately, especially considering that in ancient
Egyptian art the body colour of both man and woman was ``standardised'' and these two
colours differed noticeably from each other.  In the more ancient image the dominant
position of the husband was conspicuous, while in the New Kingdom this idea was conveyed
more weakly: although the husband was still in front of his wife, he was nevertheless
occluded by her.  One may therefore speak of the conscious use by artists of the Old
Kingdom of a geometrically contradictory image, which allowed the important idea of the
epoch to be conveyed more fully and expressively without detriment to the artistic
perfection of the work.  In that case, however, the New Kingdom composition must be
assessed not as a progressive victory of ``realism'' but as a concession to formal
geometry that had a negative effect on the ideological aspect of the composition.

In medieval art one can find not a few convincing examples of this kind.  Let us consider
the fresco ``Presentation in the Temple'' from the mural painting of Mirozhsky Monastery
in Pskov (Ill.~81).  The meeting depicted in this fresco, in the temple, between Simeon
and the prophetess Anna and Mary and Joseph, who had brought the infant Christ to the
temple, is captured at the moment when Simeon took the infant in his arms, recognising
in him the future Saviour.  Since Simeon holds the incarnate deity in his arms, the
infant is depicted above the altar and under the ciborium, i.e.\ in the holiest place of
a Christian church, where only clergy are permitted.  Here Christ is as it were on his
throne.\footnote{The legitimacy of this interpretation of the image of Christ above the
  altar and under the ciborium is confirmed, for example, by the image embroidered on the
  left shoulder-piece of the small sakkos of Metropolitan Photius (Byzantium, fourteenth
  century): ``The Saviour as a boy\ldots\ standing on the throne under a baldachin with
  Constantine and Helena in attendance'' (L.~V.~Pisarskaya.  \emph{Pamyatniki
  vizantiyskogo iskusstva V--XV vekov v Gosudarstvennoy Oruzheinoy palate} [Monuments
  of Byzantine Art of the Fifth to Fifteenth Centuries in the State Armoury Chamber].
  Leningrad--Moscow, 1964, p.~30).  The symbolic connection of Simeon, Mary, and Christ
  with the throne is already present in the Khludov Psalter miniature (ninth century),
  fol.~163 verso (see: M.~V.~Shchepkina.  \emph{Miniatyury Khludovskoy Psaltyri}
  [Miniatures of the Khludov Psalter].  Moscow, 1977).}  In order not to constrain the
movement of the figures and to show this solemn moment worthily and in all its details,
the Old Russian artist exploited the possibilities afforded him by the ambiguity of the
mapping of external space onto the picture plane.  He ``shifted'' the ciborium with its
columns backward, but placed the hands of Simeon with the infant so that visually Christ
is situated not only above the altar but also under the ciborium.

This device does not speak of the medieval artist's ``inability.''  Indeed, in the fresco
``Presentation'' from the Church of St~Panteleimon in Nerezi (Yugoslavia, twelfth
century), where Mary with Anna and Simeon with Joseph are shown at the moment when they
are still moving toward the altar and have not yet reached it, the columns of the ciborium
stand ``correctly,'' with one of them in front of the altar, partially occluding
it.\footnote{Petar Miljkovic-Pepek.  \emph{Nerezi}.  Beograd, 1966, tab.~12.}

\srcpage{181}

\figblock{81}{``Presentation in the Temple.''  Fresco of the Saviour-Transfiguration
  Cathedral of Mirozhsky Monastery, Pskov.  Twelfth century.}

On the icon ``Presentation'' from the festal tier of the iconostasis of Trinity Cathedral
of the Trinity-Sergius Lavra (Ill.~82), painted by a master of Andrei Rublev's
workshop,\footnote{Some art historians consider this icon to be by Rublev's own hand.}
the same subject is treated.  This time the artist found it possible to ``shift'' along
the line of sight not only the ciborium but the altar itself---probably in order to place
Simeon on a two-step footrest and not create an ``obstacle'' between Simeon and Mary.  In
this composition the infant Christ is also visually situated above the altar and under
the ciborium, i.e.\ the ability of the medieval artist and viewer to understand the
language of conventional shifts along the line of sight allows the deity to be connected
with the only place worthy of him---the space above the altar.  Assistance to this
perception is given by the bright vermilion-red colour of the altar, the right front
corner of which is shifted toward Simeon, and the truncated depiction of the ciborium
columns; had the artist not done the latter, the right front column would have
``screened'' the altar from Simeon, Mary, and Christ.  As for the truncation of the
columns itself, it is entirely analogous to the ``disappearance'' of the column (see
Ills.~28, 79).

\figblock{82}{Andrei Rublev, Daniil Chyorny, and masters of their circle.
  ``Presentation in the Temple.''  Icon from Trinity Cathedral of the
  Trinity-Sergius Lavra.  1420s.}

\srcpage{183}

\figblock{83}{Giotto.  ``Presentation in the Temple.''  Fresco of the Arena Chapel
  (Cappella degli Scrovegni), Padua.  C.~1305--1307.}

Disregarding the ``rules of geometry'' allows one to increase expressiveness and create
symbolic, solemn, and tranquil compositions.

Let us compare the works examined with other outstanding artistic achievements on the
same theme.  Let us cite two works by Giotto, that remarkable master of the
proto-Renaissance, who showed the same scene but strictly observed the geometric
self-consistency of the image.  In one case (Ill.~83) the artist depicted the scene of
the meeting against a background of altar with ciborium; here geometric plausibility led
to the result that the infant Christ is no longer directly connected with the space above
the altar.  He is clearly shown before the altar---this impression is reinforced by the
diagonal placement of the altar and the thick column partially concealing it.  The
everyday character of the scene intensified, and feeling this, Giotto attempts to weaken
this undesirable impression by depicting an angel.  But this novelty of Giotto's did not
gain further support; the angel could not restore to the scene the deep symbolism
inherent in it in medieval art.

\srcpage{184}

\figblock{84}{Giotto.  ``Presentation in the Temple.''  Altarpiece.  Early
  fourteenth century.  Boston, Gardner Museum.}

In the other case (Ill.~84), Giotto solves the same problem differently: he shows Christ
above the altar and under the ciborium, thereby preserving the required symbolism, but
the desire to follow geometry strictly has led to the movements of Simeon and Mary being
constrained and their figures squeezed between thin columns.  The ``realism'' of the
geometry paradoxically led to the unnaturalness of the entire scene, in which such a
strange manner of passing the infant from hand to hand is in no way justified.  It is
not surprising that this composition of Giotto's received no further development and
ceased to be used by artists.  The first of the Giotto compositions examined, in which
the infant ends up before the altar, allowed the poses of Mary and Simeon to be given a
more natural and solemn character.  Probably for this reason it and other compositions
that did not connect the infant Christ directly with the space above the altar gained
wider currency subsequently.

The ``losses'' that flow from the striving for geometric self-consistency are obvious.
All artists who painted the ``Presentation'' icons in the new manner (seventeenth to
nineteenth centuries) always depicted Simeon and the infant alongside the altar or far
from it and, in essence, transformed the deeply symbolic scene examined above, in which
its uniqueness was emphasised, into a simple illustration of the rite of
``churchings''\footnote{See, for example, the icon by Karp Zolotaryov (seventeenth
  century): Yu.~I.~Arenkova, G.~Ya.~Mekhova.  \emph{Donskoy monastyr'}.  Moscow, 1970,
  ill.~27.}---performed in Orthodox churches over every infant.  Thus illusionistic
accuracy in rendering the visible geometry of space can become an obstacle to the
conveyance of important ideas by artistic means.

Naturally, the conventional manner of painting that arises from the deliberate violation
of the geometric self-consistency of the depicted space requires appropriate preparation
in the viewer, but the latter requirement is common to all art.  This condition holds
for linear perspective too, whose ``unnaturalness'' for nearby regions of space we simply
do not notice by force of habit (and this is the appropriate preparation of the viewer);
it holds for ballet, opera, and virtually all other forms of art, and can therefore in
the slightest degree discredit the convention of Old Russian art described here.

\srcpage{185}

A completely different type of geometrically contradictory image---one that might perhaps
be called ``volumetric''---is found in Rublev's ``Trinity'' (Ill.~33).  Consider, for
example, the image of the right angel (as seen by the viewer).  If one carries out a
geometric analysis of the depiction of the throne and the angel without attention to
other objects, it can be considered self-consistent---the left legs of the throne are
covered by the depiction of the feet and garments of the angel and therefore not visible.
If one similarly examines separately the figure of the angel and the footrest, no
contradiction can be discovered there either.  However, it suffices to look attentively
at the arrangement of the figure of the angel, the throne, and the footrest simultaneously
for it to become clear that the depicted space and the objects shown in it do not accord
with each other---the throne and the footrests should intersect in space, individual parts
of them should simultaneously occupy the same points of three-dimensional objective
space, which is inconceivable for material bodies.  Geometrically there is nothing
special in this, since geometry deals not with material bodies and studies intersections
of solid figures.  Moreover, it is not excluded that the points of space simultaneously
occupied by parts of the footrest and the throne are also occupied by the right lower
part of the altar.  Literally the same can be said of the depiction of the left angel.
Thus the matter here is not shifts but another kind of geometric contradiction---the
belonging of different material points (belonging to different objects) to one and the
same point of real space.

\srcpage{186}

\figblock{85}{Two possible variants of the depiction of the throne behind the right
  angel in Andrei Rublev's icon ``Trinity.''  (Variant~$A$: the upper lateral edge of
  the throne is parallel to the nearest edge of the throne-seat; Variant~$B$: the throne
  is depicted in a natural, very mild reverse perspective.)}

If the usual ambiguity discussed above reduced to the fact that points belonging to
different objects and situated at different points of real space could correspond to one
and the same point in the picture (the objects would occlude each other), now these
points have come to coincide not only in the picture but also in objective space---which
is conceivable geometrically but physically impossible.

\srcpage{187}

The reason why such a depiction does not immediately evoke in the viewer a feeling of
inner protest and is not even noticed by him is analogous to the reason cited repeatedly
above: the visual perception of a person ``cuts off'' superfluous information once his
perceptual system has formed an unambiguous perceptual hypothesis.  An artist using such
devices must possess a great sense of tact; he must be able to ``distort'' spatial images
in such a way as to strengthen the artistic aspect of the work without overstepping the
limit permitted by the psychology of visual perception.  How fully Rublev mastered all
this is visible from the fact that after him no one has succeeded in improving the
``Trinity.''

The use of the device of ``volumetric'' geometric contradiction allows one to construct
the composition more freely and simultaneously creates a feeling of a certain
indefiniteness---the feeling that not all points have been placed over the~$i$---and
this incompleteness is probably an important feature of every outstanding work of art.
Not venturing to analyse this problem from the standpoint of art history, the author
would like to draw attention to the fact that Rublev consciously leaves other geometric
characteristics of his work indeterminate as well.  One of these questions left open is
the form of the lateral sides of the altar.  That the altar in real space has the form
of a rectangular parallelepiped is obvious; however, it remains unclear how the lateral
sides of this parallelepiped are ``depicted'' behind the angels, and yet geometrically
this is important.  Different variants are conceivable here; let us cite just two of them,
and such that there is no need to introduce an intersection of the volumes of the altar
with the throne-seats.  Let us return to the depiction of the right angel (as seen by
the viewer), since in his throne-seat two mutually perpendicular edges are shown, which
allows the form of the throne-seat to be approximately reconstructed (on the basis of
axonometric projection).  This is done in Ill.~85, where the two variants shown differ
only in the depiction of the lateral side of the altar.

Variant~$A$ is characterised by the upper lateral edge of the altar being parallel to the
nearest edge of the throne-seat.  In that case the altar has an appearance analogous to
that shown on the icon ``New Testament Trinity'' (Ill.~49): the angel sits parallel to
the lateral side of the altar, i.e.\ oriented in space the same as the middle angel, and
his rotation in the icon is connected with the fact that all three angels are depicted
from different viewpoints, analogously to Christ and the Father in the icon mentioned.
In this case the strong reverse perspective in the depiction of the altar does not
contradict the axonometric depiction of the throne-seat, since it arose as a result of
the ``gluing'' of two local axonometries, just as in the ``New Testament Trinity.''

Variant~$B$ gives a different position for the lateral edge of the altar: the altar is
depicted in a natural, very mild reverse perspective.\footnote{It is sometimes maintained
  that the strong reverse perspective of the altar is proven here by the displacement of
  the chalice toward its front edge (L.~F.~Zhegin.  \emph{Yazyk zhivopisnogo
  proizvedeniya (uslovnost' drevnego iskusstva)} [The Language of the Pictorial Work
  (Conventions of Ancient Art)].  Moscow, 1970, p.~48).  However, such shifts ``toward''
  and ``away from'' the viewer in medieval art have no relation to the perspective system;
  their reasons are explained in Chapter~VII.  In Rublev's ``Trinity'' the position of
  the chalice is compositionally determined---it harmoniously fills the space between the
  hand of the middle angel and the front edge of the altar.}  This too does not contradict
the axonometric depiction of the throne-seat, since the reverse perspective is expressed
weakly (no more strongly than in the other footrest), but now the angel sits before the
altar, and his rotation in the icon is connected with the actual rotation of the throne-seat
in real space, i.e.\ all three angels are oriented differently in space.

Not having depicted the lateral sides of the altar, Rublev consciously left the question
posed without an answer.  All researchers of the great Russian artist's work have noted
the many-sidedness of the content of the ``Trinity''; possibly hence the many-sidedness
of the geometry of the image---this geometry is indeterminate, ambiguous, as are the
ideas embodied in the ``Trinity.''

\srcpage{188}

Summing up the problem examined in the present chapter---the ambiguous correspondence of
points of the picture to points of space, and the resulting possibility of depicting
geometrically contradictory spatial images---one must first of all note that medieval
artists consciously and widely used this possibility to strengthen the artistic effect.
It is not surprising that the method of geometrically contradictory images is beginning
to attract the attention of the representatives of modern art.  Let us cite here only
one example.  Matisse's well-known canvas ``Family Portrait'' (1911, State Hermitage
Museum) contains a depiction of a chess board which, for compositional reasons, is shown
not as square ($8 \times 8$ squares) but as rectangular ($7 \times 11$ squares).

\srcpage{189}

Thus the conscious deviation from geometric self-consistency, used by both medieval and
Renaissance artists and by the masters of modern art, is no accident.  It reflects a
general principle: in the service of art, the artist has the right---even the duty---to
subordinate the formal requirements of geometry to the artistic tasks of the work.  The
examples examined here, drawn from the art of different epochs, demonstrate that this
principle has been known and applied across a vast range of artistic traditions.
